% !TeX root = ../sections/02_exercises.tex
\subsection{2-B-6.}
\begin{exercise}
	\(\alpha\in\mathbb{R}\) とする。
	数列 \(\{a_n\}_{n\in\mathbb{N}}\) の任意の部分列が
	\(\alpha\) に収束する部分列を含むならば，
	\[
	a_n\to\alpha
	\]
	であることを証明せよ。
\end{exercise}

\begin{background}
	この問題は，部分列を用いた収束判定である。
	
	通常，
	\[
	a_n\to\alpha
	\]
	ならば，任意の部分列も \(\alpha\) に収束する。
	ここではその逆に近い主張を示す。
	
	仮定は，任意の部分列そのものが \(\alpha\) に収束する，
	というものではない。
	任意の部分列の中に，さらに
	\(\alpha\) に収束する部分列が含まれる，
	という条件である。
	
	この主張を直接示すよりも，対偶を用いると分かりやすい。
	すなわち，
	\[
	a_n\not\to\alpha
	\]
	と仮定して，\(\alpha\) に収束する部分列を一切含まない部分列を作る。
\end{background}

\begin{strategy}
	対偶を示す。
	すなわち，
	\[
	a_n\not\to\alpha
	\]
	ならば，\(\alpha\) に収束する部分列を含まない部分列が存在することを示す。
	
	\(a_n\not\to\alpha\) であるから，
	ある \(\varepsilon_0>0\) が存在して，任意の \(N\in\mathbb{N}\) に対して，
	ある \(n\geq N\) が存在し，
	\[
	|a_n-\alpha|\geq \varepsilon_0
	\]
	となる。
	
	この条件を用いて，帰納的に部分列を作る。
	まず \(N=1\) として，ある \(n_1\) を取り，
	\[
	|a_{n_1}-\alpha|\geq\varepsilon_0
	\]
	とする。
	
	次に \(N=n_1+1\) として，ある \(n_2>n_1\) を取り，
	\[
	|a_{n_2}-\alpha|\geq\varepsilon_0
	\]
	とする。
	
	同様にして，
	\[
	n_1<n_2<\cdots<n_k<\cdots
	\]
	かつ，すべての \(k\) に対して
	\[
	|a_{n_k}-\alpha|\geq\varepsilon_0
	\]
	を満たす部分列
	\[
	\{a_{n_k}\}_{k\in\mathbb{N}}
	\]
	を得る。
	
	この部分列の任意の部分列を取っても，その各項は常に
	\[
	|a_{n_k}-\alpha|\geq\varepsilon_0
	\]
	を満たしている。
	したがって，その部分列は \(\alpha\) に収束できない。
	
	これは，任意の部分列が \(\alpha\) に収束する部分列を含む，
	という仮定に反する。
	
	よって対偶が示されたので，
	\[
	a_n\to\alpha
	\]
	である。
\end{strategy}